\chapter[Architettura RTL]{Architettura RTL e gerarchia dei moduli}
\label{ch:arch}

\section{Organizzazione gerarchica}
Il design è organizzato per livelli, dal moltiplicatore-accumulatore elementare fino
al top-level integrato con interfaccia SPI e PSRAM. Ogni livello incapsula il
precedente e ne astrae i dettagli: il datapath validato (\code{mac\_unit},
\code{mac8}, \code{neuron\_parallel}) non viene mai modificato dai livelli di
orchestrazione superiori.

\begin{center}
\begin{tikzpicture}[font=\footnotesize,every node/.style={fnblock,minimum width=40mm},
   level distance=13mm,sibling distance=0mm]
 \node[fnblockD,minimum width=62mm](top){\code{spi\_neuron\_top} \\ {\scriptsize top-level integrato}};
 \node[fnblockT,minimum width=62mm,below=8mm of top](arb){\code{mem\_arbiter} \;/\; \code{layer\_sequencer} \\ {\scriptsize arbitraggio 3 porte + sequenza layer}};
 \node[fnblock,minimum width=62mm,below=8mm of arb](nm){\code{neuron\_memory} \\ {\scriptsize ponte memoria $\leftrightarrow$ neurone, loop neuroni}};
 \node[fnblock,minimum width=62mm,below=8mm of nm](np){\code{neuron\_parallel} \\ {\scriptsize FSM neurone: gruppi, bias, attivazione, saturazione}};
 \node[fnblockT,minimum width=62mm,below=8mm of np](m8){\code{mac8} \\ {\scriptsize \code{PARALLEL} MAC + balanced adder tree}};
 \node[fnblock,minimum width=62mm,below=8mm of m8](mu){\code{mac\_unit} \\ {\scriptsize $x\cdot w$ + estensione segno + accumulo}};
 \foreach \a/\b in {top/arb,arb/nm,nm/np,np/m8,m8/mu}
    \draw[fnarrow] (\a) -- (\b);

 % rami memoria a destra
 \node[fnblockA,minimum width=34mm,right=14mm of nm](ma){\code{int8\_memory\_access}\\{\scriptsize byte $\leftrightarrow$ word 16-bit}};
 \node[fnblockA,minimum width=34mm,below=6mm of ma](mi){\code{memory\_interface}\\{\scriptsize handshake req/ready}};
 \node[fnblockA,minimum width=34mm,below=6mm of mi](pc){\code{psram\_controller}\\{\scriptsize bus fisico PSRAM}};
 \draw[fnarrowT] (ma)--(mi); \draw[fnarrowT] (mi)--(pc);
 \draw[fnarrowT,dashed] (nm.east) -- (ma.west);

 % rami SPI a sinistra
 \node[fnblockA,minimum width=30mm,left=14mm of arb,yshift=6mm](ss){\code{spi\_slave}\\{\scriptsize layer fisico Mode 0}};
 \node[fnblockA,minimum width=30mm,below=6mm of ss](se){\code{spi\_engine}\\{\scriptsize FSM opcode + registri}};
 \draw[fnarrowT] (ss)--(se);
 \draw[fnarrowT,dashed] (se.east) -- (arb.west);
\end{tikzpicture}
\end{center}

\section{Ruolo di ciascun modulo}
\begin{tabularx}{\textwidth}{L{3.4cm}Y}
\toprule
\rowh \thd{Modulo} & \thd{Funzione} \\
\midrule
\code{mac\_unit}       & Singolo prodotto-accumulatore: $\mathrm{acc\_out}=\mathrm{acc\_in}+(x\cdot w)$, con estensione di segno del prodotto ad \code{ACC\_WIDTH}. Parametrico su \code{DATA\_WIDTH}/\code{ACC\_WIDTH}. \\
\rowa \code{mac8}      & \code{PARALLEL} istanze di \code{mac\_unit} i cui prodotti vengono sommati da un \emph{balanced binary adder tree} di profondità $\log_2(\text{PARALLEL})$; il risultato è aggiunto all'accumulatore in ingresso. \\
\code{neuron\_parallel} & FSM di un singolo neurone: elabora \code{N\_INPUTS} ingressi in gruppi di \code{PARALLEL}, accumula tra i gruppi, somma il bias, applica l'attivazione e satura a INT8. Include il guard di elaborazione su \code{N\_INPUTS \% PARALLEL} e la larghezza runtime \code{n\_inputs\_real}. \\
\rowa \code{layer}     & Istanzia \code{N\_NEURONS} neuroni \emph{in parallelo} sullo stesso vettore di ingresso; \code{busy}=OR, \code{done}=AND dei neuroni. Percorso puramente combinatorio-di-dati usato nei benchmark del datapath. \\
\code{neuron\_memory} & Integra il calcolo con la memoria: legge $X$ (condiviso) una volta, poi per ogni neurone rilegge $W$ e bias dalla RAM e riusa una singola istanza \code{neuron\_parallel} (memory-bound, un neurone per volta). Uscita \code{y\_bus} packed neuron-major. \\
\rowa \code{layer\_sequencer} & Concatena fino a \code{N\_LAYERS} esecuzioni di \code{neuron\_memory} leggendo una tabella descrittori scritta dall'host e alternando i buffer ping-pong in RAM (Fase~5). \\
\code{act\_buffer} & Buffer di attivazione globale in block RAM \code{DP16KD}, indicizzato per id di segnale (Tipo \#2). \\
\rowa \code{graph\_engine} & Motore della rete a grafo (Tipo \#2): gather da \code{act\_buffer}, riusa \code{neuron\_parallel}, scrive le uscite per id (cap.~\ref{ch:grafo}). \\
\code{int8\_memory\_access} & Converte l'interfaccia byte/INT8 (indirizzo di byte) nell'interfaccia a parola 16-bit, selezionando il byte basso/alto tramite \code{lb\_n}/\code{ub\_n} e \code{addr>>1}. \\
\rowa \code{memory\_interface} & FSM di handshake a 2 stati (IDLE/WAIT) che serializza la singola transazione verso il controller. \\
\code{psram\_controller} & Controller del bus PSRAM parallelo asincrono con \textbf{page mode} di lettura: accesso casuale a 70~ns (\code{tAA}), burst nella stessa pagina a 20~ns (\code{tAPA}) con CE\#/OE\# tenuti attivi; abilita il page mode sul chip all'avvio via registro di configurazione (cap.~\ref{ch:mem}, \S~5.5). Pilota \code{ce\_n/oe\_n/we\_n/lb\_n/ub\_n/zz\_n} e il bus dati tri-state. \\
\rowa \code{mem\_arbiter} & Arbitro a priorità fissa (B$>$C$>$A) fra tre master byte-level: \code{spi\_engine} (A), \code{neuron\_memory} (B), \code{layer\_sequencer} (C). \\
\code{spi\_slave}     & Layer fisico SPI Mode 0, MSB-first, sincronizzatore CDC a 3 stadi su SCLK/MOSI/CS\_N, shift-register e framing di CS. \\
\rowa \code{spi\_engine} & FSM di protocollo/opcode e banco registri (\code{x\_base}, \code{w\_base}, \code{bias\_addr}, base ping-pong, attivazione, larghezze runtime\ldots), con \code{STATUS.done} sticky/clear-on-read. \\
\code{spi\_neuron\_top} & Top-level: collega SPI, arbitro, sequencer, \code{neuron\_memory} e catena PSRAM; multiplexa il controllo di \code{neuron\_memory} fra sequencer e percorso diretto single-layer. \\
\bottomrule
\end{tabularx}

\vspace{6pt}
\begin{fnnote}[Modelli di simulazione]
\code{psram\_model.v} (in \code{sim/}) e \code{memory\_model.v} sono modelli
comportamentali della memoria usati nei testbench; non fanno parte del design
sintetizzabile ma riproducono la latenza reale per la verifica end-to-end.
\end{fnnote}

\section{Due percorsi di esecuzione}
Il top-level espone due modalità mutuamente esclusive verso lo stesso motore di
calcolo \code{neuron\_memory}:
\begin{itemize}
\item \textbf{Percorso single-layer / manuale}: l'host imposta le basi con
\op{SET\_BASE}, avvia con \op{START} e legge con \op{READ\_OUTPUT}. \code{spi\_engine}
pilota direttamente \code{neuron\_memory}.
\item \textbf{Percorso multi-layer}: l'host scrive la tabella descrittori e avvia con
\op{RUN\_NETWORK}; \code{layer\_sequencer} prende possesso del controllo di
\code{neuron\_memory} (mentre \code{seq\_busy} è alto) e concatena i layer.
\end{itemize}
Il multiplexer del top-level commuta le linee di controllo di \code{neuron\_memory}
in base a \code{seq\_busy}, restituendo il motore al percorso diretto a fine sequenza.
